% META: {"id": "1NSI-RES-CO-015", "chapitre": "1NSI-RESEAUX", "type_objet": "correction", "exercice_ref": "1NSI-RES-EX-003", "capacites": ["1NSI-RESEAUX-C3"], "sources_inspiration": [], "status": "approved", "capacites_codes": ["C3"], "parcours": 1, "fichier_exercice": "chapitres/1NSI-RESEAUX/exercices/1NSI-RES-EX-015.tex"}
% BEGIN-VERIFY
% reseau = {
%     "Salle1": ["Switch1"],
%     "Salle2": ["Switch1"],
%     "Switch1": ["Salle1", "Salle2", "Switch2"],
%     "Switch2": ["Switch1", "Salle3"],
%     "Salle3": ["Switch2"],
% }
% def peuvent_communiquer(reseau, depart, arrivee, visites=None):
%     if visites is None:
%         visites = set()
%     if depart == arrivee:
%         return True
%     visites.add(depart)
%     for voisin in reseau.get(depart, []):
%         if voisin not in visites:
%             if peuvent_communiquer(reseau, voisin, arrivee, visites):
%                 return True
%     return False
% assert peuvent_communiquer(reseau, "Salle1", "Salle3") is True
% reseau2 = dict(reseau)
% reseau2["SalleIsolee"] = []
% assert peuvent_communiquer(reseau2, "Salle1", "SalleIsolee") is False
% END-VERIFY
\begin{corrige}{1NSI-RES-EX-003}
\textbf{1.} \lstinline{peuvent_communiquer(reseau, "Salle1", "Salle3")} renvoie
\lstinline{True} : il existe un chemin \lstinline{Salle1 -> Switch1 -> Switch2 -> Salle3}.

\textbf{2.} Avec \lstinline{reseau["SalleIsolee"] = []}, la fonction renvoie
\lstinline{False} : aucun chemin n'existe vers une machine qui n'a aucune connexion.

\textbf{3.} Le marquage des sommets visités évite de \textbf{reparcourir indéfiniment} un
même sommet. Sans cette précaution, un cycle (comme deux câbles entre \lstinline{Switch1}
et \lstinline{Switch2}) provoquerait des appels récursifs sans fin : la fonction
repasserait continuellement entre les deux commutateurs, sans jamais terminer
(\textbf{boucle infinie} / dépassement de la profondeur de récursion).
\end{corrige}
